Cháu chào bác,

Tối qua Hiệu gửi tin nhắn cho cháu, nhưng sáng nay cháu mới đọc. Tối qua, lần đầu tiên cháu đốt bếp củi ở nhà cháu ở Melbourne, ngọn lửa đẹp lắm, than hồng và ấm áp. Bọn trẻ con nhà cháu thích lắm. Nhớ lại ngày xưa có lần nấu bánh chưng sau nhà Đội Cấn với chị Quỳnh, bao nhiêu năm rồi.

Sáng nay, Melbourne nhiều sương. Phía trường cháu dạy học cũng gần núi, nên có lẽ sương nhiều và dày hơn. Tiết trời đã vào cuối thu. Trời không mưa tầm tã như mùa đông, cảnh vật có đượm buồn, nhưng không thê lương, giống như cảm giác của cháu lúc này.

Mỗi người nhìn thấy bác ở các góc khác nhau, một nhà Toán học, một trí thức lớn, một con người dám nói trong những hoàn cảnh không dễ dàng. Với riêng cháu, bác là con người của gia đình, của quê hương Hà Tĩnh, là người chặt gà Tết cho bác Hương, là người làm thơ viết văn và say mê sách. 

Cháu chào tạm biệt bác, để bác lại đi tiếp một chặng đường mới trong cuộc hành trình, cháu nghĩ thế, vì có ai biết được đâu là tận cùng.